\chapter{有限群的表示}\label{ch:b3:representations}

为了理解一个抽象群，让它作用于一个向量空间，
并在此作用上运用线性代数 --- 特征值、迹、内积 ---
这一纲领即\emph{表示论}，
对 $\C$ 上的有限群极为有效：每个
表示都可分解为不可约表示的直和（马施克），
不可约表示由其\emph{特征标}
（迹）唯一确定，且特征标满足正交关系，
使计算程序化。本章建立这套演算
及其最早的经典成果 --- 小群的
特征标表 --- 周末问题则导出一个远超纯群论能力范围的定理，
即伯恩赛德的 $p^aq^b$ 定理。
全文中，$G$ 为有限群，所有向量空间
均为 $\C$ 上的有限维向量空间。

\section{表示、马施克与舒尔}

\begin{definition}\label{def:b3:representations:rep}
称\emph{表示}\index{表示}为群 $G$ 的一个
同态 $\rho \colon G \to GL(V)$，这里 $\C$ 上的向量空间 $V$；
$\dim V$ 称为其\emph{次数}。子空间 $W \subseteq V$ 称为
\emph{不变}的，若 $\rho(g)W \subseteq W$ 对所有 $g$ 成立；
限制后 $W$ 成为一个\emph{子表示}。$\rho$ 称为
\emph{不可约}\index{不可约}的，若 $V \neq
0$，且其仅有的不变子空间为 $0$ 与 $V$。
\emph{同态}是 $(\rho, V)$ 与 $(\sigma, W)$ 之间的线性
映射 $f\colon V \to W$，使 $f\rho(g) = \sigma(g)f$ 对所有 $g$
成立（此即等变性）；双射 $f$ 称为\emph{同构}。
\end{definition}

\begin{example}\label{ex:b3:representations:examples}
(a) 次数为 $1$ 的表示，即 $G \to \C^\times$ 的同态。
(b) \emph{正则表示}\index{正则
表示}：$V = \C^G$，以 $(e_h)_{h \in G}$ 为基，
$\rho(g)e_h = e_{gh}$；次数为 $\abs G$。
(c) $G$ 在有限集 $X$ 上的置换作用给出
\emph{置换表示}，即 $\C^X$ 上的表示：$\rho(g)e_x =
e_{g\cdot x}$。
(d) $S_n$ 通过置换坐标作用于 $\C^n$；超平面
$\{\sum x_i = 0\}$ 是不变的，即\emph{标准
表示}，次数为 $n - 1$。
\end{example}

\begin{theorem}[马施克]\label{thm:b3:representations:maschke}
设 $W$ 为表示 $(V, \rho)$ 的任一不变子空间，
则其存在不变补空间。由此，每个表示
都是不可约表示的直和
（\emph{半单性}）。\index{马施克定理}
\end{theorem}

\begin{proof}
设 $p \colon V \to V$ 为\emph{任一}以 $W$ 为像的投影
（任取一补空间）。对其在群上取平均：
\[
\tilde p = \frac1{\abs G}\sum_{g \in G} \rho(g)\,p\,\rho(g)^{-1}.
\]
每一项将 $V$ 映入 $W$（$W$ 不变），且固定 $W$
的每一点：对 $w \in W$，$\rho(g)^{-1}w \in W$，$p$ 固定它，
而 $\rho(g)$ 将其送回原处 --- 故 $\tilde p$ 仍为
到 $W$ 的投影。它是等变的：对 $h \in G$，
$\rho(h)\tilde p\rho(h)^{-1}$ 只是对同一和式重新标号。因此 $\ker
\tilde p$ 为 $W$ 的一个不变补空间。对
各直和项迭代（因为有限维），将 $V$ 分解为不可约表示。
\end{proof}

\begin{theorem}[舒尔引理]\label{thm:b3:representations:schur}
设 $f \colon V \to W$ 为\emph{不可约}
表示之间的同态。则 $f = 0$ 或 $f$ 为同构；若
$(V,\rho) = (W,\sigma)$，则 $f = \lambda\,\mathrm{id}$，
其中 $\lambda \in \C$。由此
$\dim\operatorname{Hom}_G(V, W)$ 为 $1$（若 $V \cong W$），
否则为 $0$。\index{舒尔引理}
\end{theorem}

\begin{proof}
$\ker f$ 与 $\operatorname{im} f$ 均为不变的（由等变性），
故它们各为 $0$ 或全部：或 $f = 0$，或 $f$ 为单射
且像为全部。若 $V = W$，则 $f$ 有特征值 $\lambda$（$\C$
代数闭）；$f - \lambda\,\mathrm{id}$ 为一个
非单射的同态 $V \to V$，故为 $0$。关于
$V \cong W$ 时的维数计数：取定一个同构 $u$，任一同态
$f$ 给出自同态 $u^{-1}f = \lambda\,\mathrm{id}$：$f =
\lambda u$。
\end{proof}

\section{特征标与正交性}

\begin{definition}\label{def:b3:representations:character}
\emph{特征标}\index{特征标}为表示 $(V, \rho)$ 的
函数 $\chi_\rho(g) = \operatorname{tr}\rho(g)$。它满足
$\chi_\rho(e) = \dim V$, $\chi_\rho(hgh^{-1}) = \chi_\rho(g)$
（迹在共轭下不变）：特征标是\emph{类
函数}\index{类函数} --- 它们是空间
$\mathcal{CF}(G)$ 中的元素，即在共轭类上为常值的函数，
赋予埃尔米特内积
\[
\langle \varphi, \psi\rangle = \frac1{\abs G}\sum_{g \in G}
\overline{\varphi(g)}\,\psi(g).
\]
\end{definition}

\begin{proposition}\label{prop:b3:representations:charbasics}
$\rho(g)$ 可对角化，其特征值为单位根；
$\chi_\rho(g^{-1}) = \overline{\chi_\rho(g)}$，且
$\abs{\chi_\rho(g)} \leq \chi_\rho(e)$，且等号成立当且仅当
$\rho(g)$ 为数乘变换。特征标在直和上可加：$\chi_{V
\oplus W} = \chi_V + \chi_W$。
\end{proposition}

\begin{proof}
$\rho(g)^N = \mathrm{id}$，其中 $N = \abs G$（由拉格朗日定理）：$\rho(g)$
零化 $X^N - 1$，该多项式以单根完全分裂，故
可对角化，其特征值为 $\lambda_i \in \mu_N$。于是
$\chi(g^{-1}) = \sum \lambda_i^{-1} = \sum\bar\lambda_i =
\overline{\chi(g)}$，且 $\abs{\chi(g)} = \abs{\sum\lambda_i}
\leq \dim V$，三角不等式取等号当且仅当
所有 $\lambda_i$ 相等，即 $\rho(g) = \lambda\,\mathrm{id}$。
直和上的可加性：分块求迹。
\end{proof}

\begin{lemma}\label{lem:b3:representations:homtrace}
设 $(V, \rho)$、$(W, \sigma)$ 为表示。平均
算子定义在 $\operatorname{Hom}(W, V)$ 上，
\[
c(A) = \frac1{\abs G}\sum_{g}\rho(g)\,A\,\sigma(g)^{-1},
\]
是到 $\operatorname{Hom}_G(W, V)$ 上的投影，且对
映射 $\Phi_g \colon A \mapsto \rho(g)A\sigma(g)^{-1}$，有
$\operatorname{tr}\Phi_g =
\chi_\rho(g)\,\overline{\chi_\sigma(g)}$.
\end{lemma}

\begin{proof}
$c(A)$ 是等变的（如马施克定理证明中那样对求和作指标替换），且 $c$
保持等变映射不变（每一项都等于 $A$）：$c$ 是
以 $\operatorname{Hom}_G(W, V)$ 为像的投影。求迹：在
基下，$\Phi_g(A) = BAC$，其中 $B = \rho(g)$，$C = \sigma(g)^{-1}$；
在矩阵的基 $(E_{kl})$ 上，$BE_{kl}C = \sum_{m,n}
b_{mk}c_{ln}E_{mn}$，故 $E_{kl}$ 在
$\Phi_g(E_{kl})$ 中的系数为 $b_{kk}c_{ll}$：$\operatorname{tr}\Phi_g =
\sum_{k,l}b_{kk}c_{ll} = \operatorname{tr}(B)
\operatorname{tr}(C) = \chi_\rho(g)\chi_\sigma(g^{-1})$，且
$\chi_\sigma(g^{-1}) = \overline{\chi_\sigma(g)}$.
\end{proof}

\begin{theorem}[第一正交关系]
\label{thm:b3:representations:orthogonality}
设 $\rho, \sigma$ 为\emph{不可约}的。则\index{特征标的
正交性}
\[
\langle\chi_\sigma, \chi_\rho\rangle =
\begin{cases}
1 & \text{若 } \rho \cong \sigma,\\
0 & \text{否则：}
\end{cases}
\]
不可约特征标在
$\mathcal{CF}(G)$ 中构成标准正交族。
\end{theorem}

\begin{proof}
投影的迹等于其像的维数：
\[
\dim\operatorname{Hom}_G(W, V) = \operatorname{tr} c
= \frac1{\abs G}\sum_g \operatorname{tr}\Phi_g
= \frac1{\abs G}\sum_g \chi_\rho(g)\overline{\chi_\sigma(g)}
= \langle \chi_\sigma, \chi_\rho\rangle,
\]
而舒尔引理给出左边的值为 $\delta_{\rho \cong
\sigma}$.
\end{proof}

\begin{corollary}\label{cor:b3:representations:multiplicity}
将 $V \cong \bigoplus_i V_i^{\oplus m_i}$ 分解为互不同构的
不可约表示（$V_i \not\cong V_j$）。则 $m_i = \langle
\chi_{V_i}, \chi_V\rangle$：重数 --- 因而
在同构意义下的表示 --- 由
特征标唯一确定。此外 $\langle \chi_V, \chi_V\rangle = \sum_i
m_i^2$；特别地，$V$ 不可约当且仅当 $\langle\chi_V,
\chi_V\rangle = 1$.
\end{corollary}

\begin{proof}
$\chi_V = \sum m_i\chi_{V_i}$
（由\cref{prop:b3:representations:charbasics}）；分别与
每个 $\chi_{V_i}$ 作内积，并利用正交性。两个
特征标相等的表示有相同的重数，
从而同构。
\end{proof}

\begin{theorem}[正则表示]
\label{thm:b3:representations:regular}
设 $\chi_1, \dots, \chi_r$ 为互不相同的不可约
特征标，次数为 $n_i = \chi_i(e)$。正则
表示分解时重数为 $m_i = n_i$；
因此
\[
\sum_{i=1}^{r} n_i^2 = \abs G,
\qquad
\sum_i n_i\chi_i(g) = 0 \quad (g \neq e).
\]
\end{theorem}

\begin{proof}
正则特征标：$\chi_{\mathrm{reg}}(g) =
\#\{h : gh = h\}$，其值为 $\abs G$（$g = e$ 时），否则为 $0$。
故 $m_i = \langle \chi_i,\chi_{\mathrm{reg}}\rangle =
\frac1{\abs G}\overline{\chi_i(e)}\,\abs G = n_i$。将
$\chi_{\mathrm{reg}} = \sum n_i\chi_i$ 在 $e$ 与 $g \neq e$ 处求值，
即得上面所示的两个等式。
\end{proof}

\begin{theorem}\label{thm:b3:representations:numberirr}
不可约特征标在
$\mathcal{CF}(G)$ 中构成标准正交\emph{基}：不可约表示的个数
等于 $G$ 的共轭类个数。
\end{theorem}

\begin{proof}
只需证明完备性：设 $f \in \mathcal{CF}(G)$
与每个 $\chi_i$ 正交，证明 $f = 0$。对一个
表示 $(V,\rho)$，令 $T_{f,\rho} = \frac{1}{\abs G}
\sum_g \overline{f(g)}\,\rho(g)$。它是等变的：对 $h \in
G$,
\[
\rho(h)T_{f,\rho}\rho(h)^{-1} = \frac1{\abs G}\sum_g
\overline{f(g)}\rho(hgh^{-1})
= \frac1{\abs G}\sum_{g'}\overline{f(h^{-1}g'h)}\rho(g') =
T_{f,\rho}
\]
（$f$ 为类函数）。若 $\rho$ 为不可约表示，次数为
$n$，由舒尔引理得 $T_{f,\rho} = \lambda\,\mathrm{id}$，其中
\[
\lambda = \frac{\operatorname{tr}T_{f,\rho}}{n}
= \frac{1}{n\abs G}\sum_g \overline{f(g)}\,\chi_\rho(g)
= \frac1n\,\langle f, \chi_\rho\rangle = 0 .
\]
于是 $T_{f,\rho} = 0$ 在每个不可约表示上成立，从而（由直和）在
\emph{任一}表示 --- 特别地，在正则表示上亦成立。
将其作用于基向量 $e_e$：$0 = T_{f,\mathrm{reg}}e_e =
\frac1{\abs G}\sum_g\overline{f(g)}e_g$，从而迫使每个
$\overline{f(g)} = 0$。因此标准正交族
$(\chi_i)$ 张成 $\mathcal{CF}(G)$，其维数即为
共轭类的个数。
\end{proof}

\begin{corollary}[列正交性]
\label{cor:b3:representations:column}
对 $g, h \in G$：
\[
\sum_{i=1}^{r}\overline{\chi_i(g)}\,\chi_i(h) =
\begin{cases}
\abs{Z_G(g)} & \text{若 } g, h \text{ 共轭},\\
0 & \text{否则.}
\end{cases}
\]
\end{corollary}

\begin{proof}
设 $g_1, \dots, g_r$ 为各共轭类的代表元，$c_j$
为各类的大小。$r \times r$ 矩阵 $U_{ij} =
\sqrt{c_j/\abs G}\;\chi_i(g_j)$ 的各行标准正交
（由\cref{thm:b3:representations:orthogonality}按类写出：
$\sum_j \frac{c_j}{\abs G}\chi_i(g_j)\overline{\chi_{i'}(g_j)} =
\delta_{ii'}$），即 $UU^* = I$；方阵若满足 $UU^* =
I$ 亦满足 $U^*U = I$：各列标准正交，展开即得
上面所示的恒等式（$\abs G/c_j = \abs{Z_G(g_j)}$，即
共轭作用下的轨道--稳定子关系）。
\end{proof}

\begin{proposition}[一维特征标；提升]
\label{prop:b3:representations:onedim}
(a) $G$ 为阿贝尔群当且仅当其所有不可约表示
的次数均为 $1$；次数为 $1$ 的特征标在任意 $G$ 中的个数为
$[G : D(G)]$（它们恰为阿贝尔化的特征标）。
(b) 若 $N \trianglelefteq G$，则 $G/N$ 的不可约表示
（与 $G \to G/N$ 复合）恰好提升为 $G$ 的不可约表示中
核包含 $N$ 的那些。
\end{proposition}

\begin{proof}
(a) 若 $G$ 为阿贝尔群，则每个共轭类都是单点集：$r = \abs G$，
由 $\sum n_i^2 = \abs G$ 可得所有 $n_i = 1$；反之若
所有 $n_i = 1$，则正则表示是若干个
一维表示之和，故 $\rho_{\mathrm{reg}}(G)$ 是
同时可对角化的，从而可交换，
且 $\rho_{\mathrm{reg}}$ 忠实，故 $G$ 为阿贝尔群。次数为 $1$ 的
表示即态射 $G \to \C^\times$，其目标为阿贝尔群，
它们经由 $G^{\mathrm{ab}} = G/D(G)$ 分解，
（见\cref{exo:b3:groups:9}），
阿贝尔群 $G^{\mathrm{ab}}$ 的不同特征标个数为 $\abs{G^{\mathrm{ab}}}$（它恰有
这么多个共轭类，次数均为 $1$）。(b) 与投影复合
保持不可约性（不变子空间一一对应），
而在 $N$ 上平凡的表示
经由商群分解（见\cref{thm:b3:groups:firstiso}）。
\end{proof}

\section{特征标表}

\begin{definition}\label{def:b3:representations:table}
称\emph{特征标表}\index{特征标表}为 $G$ 的
$r \times r$ 矩阵 $\bigl(\chi_i(g_j)\bigr)$，行由不可约特征标标记，
列由共轭类标记（并标出各类的大小）。
行关于加权内积标准正交，
列正交（见\cref{cor:b3:representations:column}）：
该表严重超定，这正是它
可被计算的原因。
\end{definition}

\begin{example}[$S_3$ 的特征标表]\label{ex:b3:representations:s3}
共轭类为 $e$（1 个）、对换（3 个）、$3$-轮换（2 个）；因此 $r
= 3$ 个不可约表示，次数 $n_i$ 满足 $\sum n_i^2 = 6$：$1, 1,
2$。次数为 $1$ 的为平凡表示 $\mathbf 1$ 与符号表示 $\varepsilon$。
最后一行可由列正交性（或由
$\chi_{\mathrm{std}} = \chi_{\mathrm{perm}} - \mathbf 1$）得出：
\[
\begin{array}{c|ccc}
S_3 & e & (1\,2)\ [3] & (1\,2\,3)\ [2]\\
\hline
\mathbf 1 & 1 & 1 & 1\\
\varepsilon & 1 & -1 & 1\\
\chi_{\mathrm{std}} & 2 & 0 & -1
\end{array}
\]
验证：$\langle\chi_{\mathrm{std}},\chi_{\mathrm{std}}\rangle =
\frac{1}{6}(4 + 0 + 2) = 1$，故不可约。
\end{example}

\begin{example}[$S_4$ 的特征标表]\label{ex:b3:representations:s4}
共轭类：$e$ [1]、对换 [6]、双对换 [3]，
$3$-轮换 [8]、$4$-轮换 [6]：共 5 个不可约表示，$\sum n_i^2 =
24$，其中次数为 $1$ 的有两个（$\mathbf 1, \varepsilon$；$[S_4 :
D(S_4)] = [S_4 : A_4] = 2$）：次数分别为 $1, 1, 2, 3, 3$。
次数为 $2$ 的表示由 $S_4/V \cong S_3$ 提升而来
（见\cref{prop:b3:representations:onedim}(b)，$V$ 为克莱因四元群）；
次数为 $3$ 的表示是标准表示及其
经 $\varepsilon$ 扭转所得者：
\[
\begin{array}{c|ccccc}
S_4 & e\,[1] & (1\,2)\,[6] & (1\,2)(3\,4)\,[3] &
(1\,2\,3)\,[8] & (1\,2\,3\,4)\,[6]\\
\hline
\mathbf 1 & 1 & 1 & 1 & 1 & 1\\
\varepsilon & 1 & -1 & 1 & 1 & -1\\
\chi_2 & 2 & 0 & 2 & -1 & 0\\
\chi_{\mathrm{std}} & 3 & 1 & -1 & 0 & -1\\
\varepsilon\chi_{\mathrm{std}} & 3 & -1 & -1 & 0 & 1
\end{array}
\]
（$\chi_{\mathrm{std}}(g) = \operatorname{fix}(g) - 1$；
$\chi_2$ 是将 $S_3$ 表在各类的像上取值所得
（模 $V$）。所有行与列的正交性检验均通过 ---
任取其中两项验证即是 \cref{exo:b3:representations:3} 的热身练习。
\end{example}

\begin{method}\label{met:b3:representations:table}
构造特征标表的步骤如下：(1) 列出共轭类及其大小；
(2) 数出次数为 $1$ 的特征标，并借助 $G/D(G)$ 写出；(3)
由 $\sum n_i^2 = \abs G$ 求出其余的次数（小
整数组合问题）；(4) 易得的不可约表示：由商群提升，
由置换特征标减去 $\mathbf 1$，
（验证 $\langle\chi,\chi\rangle = 1$），以已知特征标乘以
次数为 $1$ 的特征标；(5) 借助列
正交性补全未知的行 --- 每一列与已完成的列正交，
而列 $e$ 给出各次数。
最后以完整的正交性检验核对全部结果。
\end{method}

\begin{omfigure}
\begin{tikzpicture}[scale=0.52]
  \draw[thick] (0,0) rectangle (6,6);
  \fill[omDef!25] (0,5) rectangle (1,6);
  \fill[omProp!25] (1,4) rectangle (2,5);
  \fill[omThm!20] (2,0) rectangle (6,4);
  \draw (0,5) rectangle (1,6);
  \draw (1,4) rectangle (2,5);
  \draw (2,0) rectangle (6,4);
  \node[font=\small] at (0.5,5.5) {$1$};
  \node[font=\small] at (1.5,4.5) {$1$};
  \node[font=\small] at (4,2) {$M_2(\C)$};
\end{tikzpicture}

{\small $S_3$ 的正则表示，作块对角化：
$\C[S_3] \cong \C \times \C \times M_2(\C)$，维数 $1 + 1 +
4 = 6 = \abs{S_3}$。一般地 $\C[G] \cong \prod_i
M_{n_i}(\C)$：恒等式 $\sum n_i^2 = \abs G$
是关于矩阵块的陈述。}
\end{omfigure}

\section{练习}

\begin{exercise}[$\star$]\label{exo:b3:representations:1}
(a) 证明 $\Z/n\Z$ 的不可约特征标为
$\chi_k(\bar m) = \eu^{2\iu\pi km/n}$，$k = 0, \dots, n-1$，并
写出 $\Z/4\Z$ 的特征标表。
(b) 验证其上的两个正交关系 --- 并认出
该矩阵：本书何处见过它？
\end{exercise}

\begin{exercise}[$\star$]\label{exo:b3:representations:2}
设 $G$ 作用于有限集 $X$，记 $\chi$ 为
置换表示 $\C^X$ 的特征标。
(a) 证明 $\chi(g) = \abs{\operatorname{Fix}_X(g)}$，并
$\langle \mathbf 1, \chi\rangle = \#\{\text{轨道数}\}$ ---
伯恩赛德计数引理（\cref{exo:b3:groups:5}）
是一次特征标计算。
(b) 设该作用可迁，从而 $\chi = \mathbf 1 +
\psi$。证明 $\psi$ 不可约当且仅当该作用是
\emph{$2$-可迁}（在互异点的有序对上
可迁）。\emph{（将 $\langle\chi,\chi\rangle$ 计为
$X \times X$ 上的轨道数。）}
(c) 由此得出 $S_n$（$n \geq
2$）的标准表示不可约。
\end{exercise}

\begin{exercise}[$\star$]\label{exo:b3:representations:3}
从头重建 $S_3$ 的特征标表，依
\cref{met:b3:representations:table}，然后验证
两行与两列的正交关系，于 $S_4$ 的特征标表
（见\cref{ex:b3:representations:s4}）。
分解 $S_4$ 作用于 $\{1,2,3,4\}$ 的置换特征标，以及
$\chi_{\mathrm{std}}^2$（逐点平方）分解为不可约特征标。
\end{exercise}

\begin{exercise}[$\star\star$]\label{exo:b3:representations:4}
计算 $D_4$ 与 $Q_8$ 的特征标表，由此得出
两个不同构的群可以具有相同的特征标表 ---
这张表在此对群中仍刻画了哪些群论数据
（中心的阶、阿贝尔化、
对合元个数）？其中哪些它\emph{未能}
刻画？
\end{exercise}

\begin{exercise}[$\star\star$]\label{exo:b3:representations:5}
$A_4$ 的特征标表：共轭类 $e$ [1]、双对换
[3]，以及\emph{两}个 $3$-轮换共轭类 [4]、[4]。
(a) 解释 $3$-轮换的分裂（比较
$S_4$ 与 $A_4$ 中的中心化子，参照\cref{exo:b3:groups:11}）。
(b) 求三个次数为 $1$ 的特征标（借助 $A_4/V \cong
\Z/3\Z$），以及次数为 $3$ 的特征标（限制
$\chi_{\mathrm{std}}$ 自 $S_4$），并拼出该表。
(c) 由该表读出 $A_4$ 的正规子群
（核 $\{g : \chi(g) = \chi(e)\}$ 及其交集）。
\end{exercise}

\begin{exercise}[$\star\star$]\label{exo:b3:representations:6}
(a) 证明 $\ker\chi = \{g : \chi(g) = \chi(e)\}$ 是
所对应表示的核
（\cref{prop:b3:representations:charbasics}，取等情形）。
(b) 证明 $G$ 的每个正规子群都是
不可约特征标之核的交。\emph{（忠实表示 $G/N$：
取其正则表示。）}
(c) 由此推出：$G$ 为单群，当且仅当 $\ker\chi_i = \{e\}$
对每个非平凡不可约 $\chi_i$ 成立 --- 单性可从
特征标表中读出。
\end{exercise}

\begin{exercise}[$\star\star$]\label{exo:b3:representations:7}
$\abs G = 8$ 时，证明阶为 $8$ 的非交换群
次数型为 $(1,1,1,1,2)$，且其次数为 $2$ 的
表示是忠实的。更一般地，证明阶为 $p^3$ 的非交换群
次数型为 $(1^{\,p^2},
p, \dots, p)$，即 $p^2$ 个 1 与 $p - 1$ 个
次数为 $p$ 的特征标。\emph{（用 $[G : D(G)]$ 与 $\sum n_i^2 = \abs G$；此处
$D(G) = Z(G)$ 的阶为 $p$。）}
\end{exercise}

\begin{exercise}[$\star\star\star$]\label{exo:b3:representations:8}
设 $G, H$ 为有限群：证明
类函数 $\chi(g)\psi(h)$（定义在 $G \times H$ 上），其中 $\chi, \psi$
分别是 $G, H$ 的不可约特征标，恰是
$G \times H$ 的全部不可约特征标。\emph{（标准正交性是直接计算；
完备性通过计数共轭类得到。）}由此推出
$\Z/2\Z \times \Z/2\Z$ 的特征标表，并借助结构定理重新推导
\cref{prop:b3:representations:onedim}(a) 在
有限交换群上的结论。
\end{exercise}

\begin{exercise}[$\star\star\star$]\label{exo:b3:representations:9}
$A_5$ 的特征标表（共轭类大小依次为 $1, 15, 20, 12,
12$，见 \cref{exo:b3:groups:11}）：
(a) 证明次数为 $1, 3, 3, 4, 5$ \emph{（唯一满足 $\sum n_i^2 = 60$
且 $n_1 = 1$ 的解，
利用 \cref{exo:b3:representations:6}(c) 与单性排除
其他 $n_i = 1$ 的情形）}。
(b) 构造次数为 $4$ 的特征标（置换作用于
$5$ 个点上）以及次数为 $5$ 的特征标（作用于六个
西罗 $5$-子群上，得次数 $6 = 1 + 5$；
需验证其不可约），再由列正交性补全两行次数为 $3$ 的
特征标，黄金分割型取值 $\frac{1\pm\sqrt5}2$ 出现在
$5$-轮换所在的共轭类上。
(c) 由已完成的特征标表验证 $A_5$ 是单群，
（见 \cref{exo:b3:representations:6}(c)）。
\end{exercise}

\begin{exercise}[$\star\star$]\label{exo:b3:representations:10}
设 $\rho$ 为次数为 $n$ 的不可约表示，且 $z
\in Z(G)$。证明 $\rho(z) = \lambda_z\,\mathrm{id}$，其中
$\lambda \colon Z(G) \to \C^\times$ 为同态（
\emph{中心特征}），由此推出 $\abs{\chi(z)} = n$，
对中心元 $z$ 成立。应用：若 $G$ 有一个忠实的不可约表示，
则 $Z(G)$ 为循环群。
\end{exercise}

\begin{exercise}[$\star\star$]\label{exo:b3:representations:11}
（同型投影）设 $(V, \rho)$ 为
$G$ 的表示，$\chi_i$ 为次数为 $n_i$ 的不可约特征标。
记
\[
p_i = \frac{n_i}{\abs G}\sum_{g\in G}
\overline{\chi_i(g)}\,\rho(g) \;\in\; \mathcal L(V).
\]
(a) 证明 $p_i$ 是 $G$-等变的，并计算它
在不可约子表示 $W \subseteq
V$（特征标为 $\chi_j$）上的限制，结果为 $\delta_{ij}\,
\mathrm{id}_W$ \emph{（由舒尔引理，
取迹确定纯量）}。
(b) 由此推出 $p_i$ 是到 $V_i$ 上的投影，
即特征标为 $\chi_i$ 的全部不可约子表示之和（
\emph{同型分支}），并且 $\sum_ip_i =
\mathrm{id}_V$，且分解 $V =
\bigoplus_iV_i$ 是典范的 --- 不同于
把每个 $V_i$ 分解为不可约表示的更细分解。
(c) 对 $S_3$ 的正则表示与符号
特征标 $\varepsilon$，将 $p_\varepsilon$ 显式写成
群代数中的元素，并直接验证 $p_\varepsilon^2
= p_\varepsilon$。
\end{exercise}

\begin{exercise}[$\star\star$]\label{exo:b3:representations:12}
（读表）某群 $G$ 的特征标表部分已知：
其阶为 $24$，共有 $5$ 个共轭类，
大小依次为 $1, 6, 8, 6, 3$，次数依次为 $1, 1, 2, 3, 3$。
(a) 求出完整的特征标表：先求两个线性特征标（一个为平凡的，
另一个恰在共轭类上取值 $-1$，
即大小为 $6$ 和 $6$ 的那两类），再求次数为 $2$ 的特征标，
由与单位元列的列正交性求出；然后
同样求两个次数为 $3$ 的特征标，其中之一为 $\chi_2\chi_4$。
(b) 确定 $G$（$\cong S_4$：比较共轭类与轮换型），
并由特征标表求出正规子群，
利用 \cref{exo:b3:representations:6}：$\chi_2$ 的核（指数
为 $2$，即 $A_4$）与 $\chi_3$ 的核（即克莱因四元群 $V_4$），除此之外只有
$\{e\}, G$。
(c) 说明特征标表如何表明 $G/V_4 \cong S_3$
\emph{（哪些特征标可经由该商群得到？）}。
\end{exercise}

\section{问题：伯恩赛德 \texorpdfstring{$p^aq^b$}{paqb}
定理}

\begin{problem}[{周末问题 --- 阶为
$p^aq^b$ 的群的可解性}]\label{pb:b3:representations:1}
伯恩赛德于 1904 年证明：阶至多有两个素因子的群都是可解的 ---
这是关于抽象群的一个结论，半个世纪以来，
唯一已知的证明都依赖特征标理论。
本题将完整构造这一证明，综合运用 \cref{ch:b3:groups}（可解性）、
\cref{ch:b3:modules}（有限生成 $\Z$-模）以及本章的内容。
下文中，$\chi_1, \dots, \chi_r$ 为 $G$ 的全部
不可约特征标，
且 $n_i = \chi_i(e)$。

\medskip
\noindent\textbf{第一部分 --- 代数整数。}
\emph{代数整数}是某个\emph{首一}多项式的根，
该多项式属于 $\Z[X]$。
\begin{enumerate}
  \item 证明 $\alpha$ 是代数整数，当且仅当环
        $\Z[\alpha]$ 是有限生成的 $\Z$-模。
  \item 由此推出全体代数整数构成 $\C$ 的
        子环。\emph{（若 $\Z[\alpha], \Z[\beta]$ 有限生成，
        则 $\Z[\alpha, \beta]$ 亦有限生成，有限生成 $\Z$-模的子模
        仍有限生成，
        由 \cref{thm:b3:modules:submodule} 及呈现论证即得；
        或直接地，$\Z^n$ 的子模是自由模，
        其秩 $\leq n$）}。
  \item 证明 \emph{有理}代数整数
        必为整数。\emph{（有理根定理。）}
  \item 证明每个特征标值 $\chi(g)$ 都是
        代数整数。
\end{enumerate}

\medskip
\noindent\textbf{第二部分 --- 类和关系。}固定
不可约表示 $(V, \rho)$，次数为 $n$，特征标为 $\chi$。对一个
共轭类 $C$，令 $S_C = \sum_{g \in C}\rho(g) \in
\mathcal L(V)$。
\begin{enumerate}[resume]
  \item 证明 $S_C$ 是等变的，从而 $S_C =
        \omega_C\,\mathrm{id}$，其中
        \[
        \omega_C = \frac{\abs C\,\chi(g_C)}{n}
        \qquad (g_C \in C \text{ 任一代表元}).
        \]
  \item 证明 $S_CS_{C'} = \sum_{C''}
        a_{CC'C''}\,S_{C''}$，其中 $a_{CC'C''} \in \N$ 为
        对固定的 $z \in C''$，满足 $(x, y) \in C \times
        C'$ 且 $xy = z$。由此推出 $\omega_C$ 满足
        $\omega_C\,\omega_{C'} = \sum_{C''}
        a_{CC'C''}\,\omega_{C''}$。
  \item 由此推出每个 $\omega_C$ 是代数整数。
        \emph{（$\Z$-模由 $1$ 与各
        $\omega_C$ 生成）是有限生成的环；应用第 1 小题的判别准则 ---
        更确切地说，证明 $M = \Z\text{-张成}
        (1, (\omega_C)_C)$ 满足 $\omega_{C_0} M \subseteq
        M$ 并利用行列式与凯莱–哈密顿技巧，或
        第 2 小题的子环论证。）}
  \item 由此推出 \emph{弗罗贝尼乌斯整除性}：$n_i$ 整除
        $\abs G$，对每个不可约次数均成立。
        \emph{（计算 $\frac{\abs G}{n} = \frac{\abs G}{n}
        \langle\chi,\chi\rangle = \sum_C \omega_C\,
        \overline{\chi(g_C)}$：这是一个代数整数，
        且为有理数。）}
\end{enumerate}

\medskip
\noindent\textbf{第三部分 --- 伯恩赛德单群判别法。}
\begin{enumerate}[resume]
  \item 设 $\chi$ 为次数为 $n$ 的不可约特征标，$C$ 为一个
        共轭类，且 $\gcd(\abs C, n) = 1$。利用贝祖定理及第
        4--7 题，证明 $\frac{\chi(g_C)}{n}$ 是一个
        代数整数。
  \item 进一步设 $0 < \abs{\chi(g_C)} < n$。证明这是不可能的：
        代数整数 $\alpha =
        \chi(g_C)/n$ 的全部\emph{共轭} --- 即
        数 $\alpha_\sigma = \frac1n\sigma(\chi(g_C))$，其中
        $\sigma \in \operatorname{Gal}(\Q(\zeta_{\abs G})/\Q)$，
        每个都是 $n$ 个单位根的平均 --- 模长
        $\leq 1$，
        从而乘积 $N = \prod_\sigma\alpha_\sigma$ 是
        有理代数整数，且 $0 < \abs N < 1$ ---
        逐条论证上述断言，引用
        \cref{thm:b3:galois:cyclotomicirred} 给出的伽罗瓦群结论，
        以及 $\sigma$ 置换单位根这一事实。
        由此得到：要么 $\chi(g_C) = 0$，要么 $\rho(g_C)$
        为数量矩阵（\cref{prop:b3:representations:charbasics}）。
  \item （伯恩赛德判别法）设 $C \neq \{e\}$ 为
        大小为 \emph{素数幂} $p^k > 1$ 的共轭类，且设 $G$ 为
        非交换单群。对 $C$ 所在的列
        与 $e$ 所在的列应用列正交性，得到
        $1 + \sum_{i \geq 2} n_i\chi_i(g_C) = 0$。证明
        存在非平凡的 $\chi_i$（$p \nmid n_i$），使得
        $\chi_i(g_C) \neq 0$ \emph{（否则 $\frac1p$ 将是
        代数整数）}。由第 10 题，$\rho_i(g_C)$
        为数量矩阵，由此推出与单性相矛盾；
        \emph{（$g$ 的集合，其中 $\rho_i(g)$ 为数量矩阵，
        构成正规子群；利用
        \cref{exo:b3:representations:6} 中的忠实性）}。结论：\emph{任一非交换单群
        都不存在大小为素数幂的共轭类，
        且大小 $> 1$。}
\end{enumerate}

\medskip
\noindent\textbf{第四部分 --- 主定理。}
\begin{enumerate}[resume]
  \item 设 $\abs G = p^aq^b$，其中 $a + b \geq 1$。若 $G$ 为
        单群，证明它是交换群：取 $z \neq e$，使其为
        某西罗 $q$-子群的中心元
        （\cref{thm:b3:groups:pfixed}），考虑其共轭类的大小
        $[G : Z_G(z)]$，它是 $p$ 的幂，
        （为什么？），再应用第 11 题。
  \item 对 $\abs G$ 作归纳，\emph{阶为 $p^aq^b$ 的
        每个群都是可解群}（伯恩赛德）。为何该论证
        在涉及三个素数时失效——是否必然如此，
        鉴于 $\abs{A_5} = 2^2\cdot3\cdot5$？
\end{enumerate}

\medskip
\noindent\textbf{第五部分 --- $A_5$ 的特征标表。}伯恩赛德定理
       无法触及的最小群值得完整刻画，
       以下内容只用到本章内容，以及
       \cref{exo:b3:groups:11}。
\begin{enumerate}[resume]
  \item 回忆 \cref{exo:b3:groups:11} 中
        $A_5$ 的五个共轭类：$\{e\}$，$15$ 个双对换，
        $20$ 个三轮换，以及两个
        各含 $12$ 个五轮换的类，分别以 $c =
        (1\,2\,3\,4\,5)$ 和 $c^2$ 为代表。试说明
        为何这些五轮换在 $A_5$ 中分裂为两个共轭类，
        而在 $S_5$ 中却同属一个共轭类？
  \item 证明 $A_5$ 的不可约次数恰为
        $1, 3, 3, 4, 5$：利用 $\sum_in_i^2 = 60$ 及 $r = 5$
        个共轭类，且 $A_5$ 是完全群（$D(A_5) =
        A_5$），故平凡特征标是其唯一的线性特征标；
        然后排除其余所有多重集\emph{（将 $59$
        写成四个整数（$\geq 2$）的平方和，验证
        这只能以一种方式成立，且所有加数都是合理的
        次数）}。
  \item 设 $\pi$ 为 $A_5$ 作用在
        $\{1, \dots, 5\}$ 上的置换特征标：$\pi(g) = \#\operatorname{Fix}(g)$，
        它在五个共轭类上的值为 $5, 1, 2, 0, 0$。计算
        $\langle\pi, \mathbf 1\rangle$ 与 $\langle\pi,
        \pi\rangle$，并由此推出 $\chi_4 = \pi - \mathbf 1$
        是次数为 $4$ 的不可约特征标，其值为 $4, 0, 1, -1,
        -1$。
  \item 对 $10$ 个无序对 $\{i, j\}$ 作同样的讨论：
        不动点个数依次为 $10, 2, 1, 0, 0$。计算
        $\langle\pi_{10}, \pi_{10}\rangle$、$\langle\pi_{10},
        \mathbf 1\rangle$ 及 $\langle\pi_{10}, \chi_4\rangle$，
        由此推出分解 $\pi_{10} = \mathbf 1 + \chi_4
        + \chi_5$，并得到不可约特征标 $\chi_5$，
        次数为 $5$，其值为 $5, 1, -1, 0, 0$。
  \item 其余两个不可约特征标 $\chi_2, \chi_3$ 次数均为
        $3$。列正交性（非单位列与单位列之间，以及每列与自身之间）
        决定了它们在五轮换之外的取值：
        证明在双对换上有
        $\chi_2(g) = \chi_3(g) = -1$，
        在三轮换上为 $0$。
  \item 在五轮换所在的共轭类上，记 $x = \chi_2(c)$，$y =
        \chi_2(c^2)$；由对称性，可取 $\chi_3(c) = y$，
        $\chi_3(c^2) = x$。由 $c$ 所在列分别与单位列及 $c^2$ 所在列
        作正交，推出
        $x + y = 1$ 与 $xy = -1$（并检验 $x^2 +
        y^2 = 3$ 为 $c$ 所在列与自身所给出的值），于是
        \[
        \{x, y\} = \Bigl\{\frac{1 + \sqrt5}2,\ \frac{1 -
        \sqrt5}2\Bigr\} :
        \]
        即黄金比及其共轭。组装
        出 $A_5$ 的完整特征标表。
  \item 进行检验：$\chi_2$ 的行范数为 $1$（利用
        $\varphi^2 + \bar\varphi^2 = 3$），$\langle\chi_2,
        \chi_3\rangle = 0$，以及弗罗贝尼乌斯整除性（第
        8 题）对全部五个次数成立。在特征标表中，何处能
        \emph{看出}它与 $S_5$ 的差异，其全部
        特征标值都是有理整数？
  \item 仅由该表即可推出 $A_5$ 是单群：正规子群是
        若干共轭类之并，
        这些类含 $e$，其元素个数整除 $60$ --- 检验
        $1 + 15 + 20 + 12 + 12$ 的真部分和中，凡包含项 $1$ 者，
        均不整除 $60$。与第 11 题的判别法
        相互印证：验证 $A_5$ 的各共轭类中，
        没有大小为素数幂且 $> 1$ 的类。
  \item （正二十面体尾声）$A_5$ 是正二十面体的旋转群，
        次数为 $3$ 的表示是
        $\R^3$ 上的两种几何作用。验证迹
        恒等式：绕角 $\theta$ 的旋转，其迹为 $1 +
        2\cos\theta$，又 $1 + 2\cos\frac{2\pi}5 =
        \frac{1+\sqrt5}2 = \varphi$。无需计算，
        即可解释为何另一个次数为 $3$ 的特征标
        必须取共轭值。事实上，$\Q(\sqrt5)/\Q$ 的伽罗瓦群
        逐项作用于整个特征标表，
        即对每个分量分别作用，置换其中的不可约特征标。
\end{enumerate}

\medskip
    \noindent\textbf{第六部分 --- 补充：一个中心界与
张量平方}
\begin{enumerate}[resume]
  \item （比整除性更强的一个结论）设 $\chi$ 为不可约
        特征标，次数为 $n$。证明 $\abs{\chi(z)} = n$ 对每个
        $z \in Z(G)$ 成立 \emph{（舒尔引理：$\rho(z)$ 为
        标量，阶有限）}，并由
        $\langle\chi, \chi\rangle = 1$ 推出如下的界
        \[
        n^2 \leq [G : Z(G)] .
        \]
        证明：$8$ 阶非交换群的不可约次数为
        $1, 1, 1, 1, 2$ \emph{（共五个
        共轭类；将 $8$ 写成五个平方数之和）}，
        并取到等号 $4 = [G : Z(G)]$；验证该界在
        $A_5$ 上成立，其中心是平凡的。
  \item （$\chi_2$ 的张量平方）对有限阶元素 $g$，
        $\rho(g)$ 可对角化，且特征值为单位根；
        由此推出特征标公式
        \[
        \chi_{\operatorname{Sym}^2 V}(g)
        = \frac{\chi(g)^2 + \chi(g^2)}2,
        \qquad
        \chi_{\Lambda^2 V}(g)
        = \frac{\chi(g)^2 - \chi(g^2)}2 .
        \]
        将上述公式用于 $\chi_2$，它是 $A_5$ 的特征标 \emph{（注意 $g^2$
        取遍 $c^2$ 所在共轭类时，$g$ 取遍 $c$ 所在共轭类，反之亦然
        ）}，证明 $\Lambda^2\chi_2 =
        \chi_2$ 及 $\operatorname{Sym}^2\chi_2 = \mathbf 1 +
        \chi_5$，于是
        \[
        \chi_2\otimes\chi_2 = \mathbf 1 + \chi_2 + \chi_5 .
        \]
        几何地解释 $\Lambda^2\chi_2 = \chi_2$，
        即借助 $\R^3$ 上的叉积。
  \item （表的最终核对）数值验证两个五阶轮换所在列之间的
        正交性，
        即 $\varphi\bar\varphi = -1$，并验证 $5 =
        \abs{Z_{A_5}(c)}$，取 $c$ 所在列与自身作内积；
        以及正则特征标 $\sum_i
        n_i\chi_i(g) = 0$ 在表中四个非单位元所在的
        列上为零。
\end{enumerate}
\end{problem}
